\section{Navigation for Source Section I}
\label{sec:introduction-navigation}

\sourceeqrange{background, motivation, and roadmap; no main equation inventory}

\subsection{Purpose of This Companion Section}

Source Section I of \OnukiPRE{} motivates a dynamic van der Waals model for a
one-component fluid undergoing gas--liquid transition in nonuniform temperature.
Rather than retelling that motivation, this companion gives a reading map:
what physical problem is being prepared, what mathematical objects the reader
should expect in Section II, and what later scenario descriptions in Section III
depend on.

The section is navigation only.  It contains no new derivation, no symbolic
identity, and no independent historical survey.  The original article remains
the authority for the source narrative and for the bibliographic context.

\subsection{Physical Setting Prepared by the Introduction}

The source paper combines four ingredients that should be kept separate while
reading the later equations.

\begin{enumerate}[label=\arabic*.]
\item \textbf{Bulk van der Waals thermodynamics.}  The local variables are the
  number density, temperature, entropy, internal energy, pressure, and chemical
  potential.  Section~\ref{sec:vdw-baseline} derives the specific formulas used
  as the baseline.
\item \textbf{Diffuse-interface capillarity.}  A square-gradient contribution is
  added to describe interfaces without imposing a sharp moving boundary.
  Sections~\ref{sec:vdw-baseline}--\ref{sec:equilibrium-conditions} explain the
  equilibrium gradient identities and wall branch.
\item \textbf{Hydrodynamic transport.}  Density, momentum, energy, entropy,
  viscous stress, heat flux, and reversible capillary stress enter as local
  balance-law objects.  Section~\ref{sec:hydrodynamic-equations} derives the
  bookkeeping that the introduction only prepares.
\item \textbf{Nonisothermal phase-change scenarios.}  The examples later in the
  source involve heat flow, evaporation or condensation, and interface motion.
  The companion treats those examples as reading guides and finite estimate
  checks, not as reproduced simulations.
\end{enumerate}

This separation is important.  A constitutive law, a balance law, a gradient
identity, and a numerical scenario have different logical status.  The companion
will not turn the introductory motivation into a proof of thermodynamic
validity, PDE well-posedness, or numerical correctness.

\begin{figure}[htbp]
\centering
\resizebox{\linewidth}{!}{\begin{tikzpicture}[
  font=\small,
  node distance=10mm and 9mm,
  box/.style={draw=black, line width=0.45pt, rounded corners=2pt,
    align=center, text width=31mm, minimum height=11mm, fill=gray!8},
  note/.style={draw=black, line width=0.35pt, rounded corners=2pt,
    align=center, text width=37mm, minimum height=10mm, fill=white},
  arrow/.style={-{Latex[length=2.2mm]}, line width=0.5pt},
  darrow/.style={-{Latex[length=2.2mm]}, dashed, line width=0.5pt}
]
\node[box] (i) {Source I\\motivation and scope};
\node[box, right=of i] (iia) {II.A\\bulk thermodynamics};
\node[box, right=of iia] (iib) {II.B--II.C\\gradient and wall branches};
\node[box, right=of iib] (iid) {II.D\\hydrodynamic balances};
\node[box, below=of iid] (app) {Appendices A--B\\stress and scaling};
\node[box, left=of app] (iii) {Section III\\scenario estimates};
\node[note, below=of iib] (boundary) {Global wall, boundary, and simulation-code branches remain separately recorded};

\draw[arrow] (i) -- (iia);
\draw[arrow] (iia) -- (iib);
\draw[arrow] (iib) -- (iid);
\draw[arrow] (iid) -- (app);
\draw[arrow] (app) -- (iii);
\draw[darrow] (iib) -- (boundary);
\draw[darrow] (iid) -- (boundary);
\draw[darrow] (app) -- (boundary);
\draw[darrow] (iii) -- (boundary);
\end{tikzpicture}
}
\caption{Source-section navigation used by the companion.  Solid arrows show
the reading order of source-supported derivations.  Dashed arrows mark
boundary, wall-transfer, and simulation-code issues that remain separate from
local equation derivations.}
\label{fig:source-section-navigation}
\end{figure}

\subsection{Reader Prerequisites and Notation Warnings}

Before entering Section II, the reader should be comfortable with the following
minimum toolkit:

\begin{itemize}
\item thermodynamic partial derivatives at fixed density, temperature, entropy,
  or volume-like variables;
\item variational derivatives of one-dimensional and multidimensional
  square-gradient functionals;
\item local conservation laws written as density-plus-divergence balances;
\item tensor notation for stress divergence and stress power;
\item the distinction between local differential identities and integrated
  balances with boundary or wall terms.
\end{itemize}

The notation changes role as the article proceeds.  For example, the coefficient
called \(M\) first appears in an equilibrium square-gradient free-energy density
and later becomes the nonequilibrium combination \(M=CT+K\).  Likewise,
pressure in the bulk thermodynamic identity should not be confused with the full
reversible pressure tensor.  These changes are tracked in the notation
crosswalk and in the section-specific symbol registries.

\subsection{Roadmap from Section I to the Companion Derivations}

The introduction points toward the following source-supported derivation order.

\begin{itemize}
\item \textbf{Section II.A:} establish the local van der Waals thermodynamic
  baseline, spinodal and sound quantities, and the equilibrium gradient
  interface formulas.
\item \textbf{Section II.B:} replace purely isothermal gradient free energy by
  entropy and energy functionals with gradient contributions, leading to
  \(M=CT+K\) and the generalized chemical-potential variation.
\item \textbf{Section II.C:} record equilibrium and wall branches, including
  the natural density boundary condition.
\item \textbf{Section II.D:} introduce hydrodynamic balances, total/internal
  energy bookkeeping, entropy-production bookkeeping, reversible stress,
  reversible entropy flux, and boundary-dependent global entropy statements.
\item \textbf{Appendices A--B:} supply the reversible-stress derivation and the
  scaled-equation branch analysis, including the printed-B3 versus
  dimensional-source distinction.
\item \textbf{Section III:} read the numerical examples as source scenarios with
  stated scales, assumptions, and estimates.  This companion does not infer the
  unpublished simulation code or reproduce the simulations.
\end{itemize}

This order is a study-guide order, not a claim that the introduction itself
contains the equations.  The equations enter in the source sections listed
above and are cited there by number.

\subsection{Boundaries of the Introduction Guide}

The introduction refers to broad capillarity and phase-field traditions.  This
companion records local context sources such as Slemrod and
Dunn--Serrin only as background for future audits
\citep{Slemrod1983VanDerWaalsPhaseBoundaries,DunnSerrin1985}.  It does not use
those sources as derivation authority until exact page, equation, variable, and
boundary anchors are recorded.  The direct Korteweg 1901 source has now been
page-image inspected, including its pressure-positive motion convention and
four-function capillary stress in Eq.~(20) \citep{Korteweg1901}.  Appendix~
\ref{app:korteweg-primary-crosswalk} keeps that source transcription separate
from a modern tensor reconstruction and from the conditional comparison with
Onuki.  This establishes neither a formula genealogy nor full model
equivalence.

The guide also keeps later source-dependent issues out of the introduction:

\begin{itemize}
\item the Appendix-B printed-B3 versus dimensional-source branch remains a later
  scaled-equation issue;
\item the unpublished numerical implementation branch remains
  \status{UNRESOLVED\_SOURCE\_DEPENDENT};
\item global entropy and energy balances remain boundary/wall-transfer
  dependent;
\item Section III scenarios remain reading guides unless a later human-approved
  numerical reproduction pass is opened.
\end{itemize}

\subsection{Verification and Open Issues}

No algebraic or dimensional software check is attached to this section.  The
verification requirement is source safety: the prose must be original, must not
replace the source introduction, and must not add unsupported historical or
simulation claims.  The relevant project checks are the bibliography/provenance
checks, the section-proof checklist, the PDF build, and human review of the
rendered page.
